\scanpage{77}
connexe, dont le $\pi_1$ est isomorphe (en tant que \uline{groupe extérieur}) :
\[
(95)\quad\pi_1(\Pi M_{\underline G})=\pi_1(M_{\underline G})\simeq\prod_{\alpha\in S(G)}\pi_1(M_{g_\alpha,\widehat I_\alpha})\simeq\prod_\alpha\mathcal T^+_{g_\alpha,I_\alpha}\simeq\prod_\alpha\mathcal T^{!+}_{g_\alpha,\nu_\alpha}.
\]
\margin{où les $\mathcal T^{!+}_{g,\nu}$ sont des groupes de Teichmüller}
Les opérations de $\unsure{\Gamma^{!}_{\underline G}|\Gamma_{\underline G}}$ sur $\Pi M_{\underline G}$ définissent un groupoïde mixte $(\Pi M_{\underline G},\Gamma_{\underline G})$, dont le groupe fondamental (extérieur) s'insère \deleted{de \unsure{par}} \inserted{dans} une structure d'extension
\[
(96)\quad1\longrightarrow\pi_1(\Pi M_{\underline G},\Gamma_{\underline G})\longrightarrow\pi_1(\Pi M_{\underline G},\Gamma_{\underline G})\unsure{-|\longrightarrow}\Gamma_G\longrightarrow1
\]
Nous allons définir une \inserted{sorte de} variance des $\Pi M_{\underline G}$ relativement à deux sortes de ``morphismes'' pour les MD-graphes $\underline G$. 

Tout d'abord une \uline{``variance de générisation''}, qui est composée de la \uline{variance pour isomorphismes} de MD-graphes (que j'admets comme ``triviale'', et qu'il est inutile de préciser) et une ``variance de passage au quotient de MD-graphes''.
\margin{76}
